% page 184 of the facsimile (image p-183 of the scan)
% block 157.5 x 233.3 mm ; left margin 8.0 mm
\pgc{-12.700mm}{0.000mm}{
\l{\cel{4}176}
\l{}
\l{\sou{fosfeno}. (\sou{fiziol.}) Fenomeno lumala, quan onu povas produktar}
\l{\cel{3}sur la retino per kompresar la okul-globo. - DEF.}
\l{}
\l{\sou{fosforecar}. (\sou{netrans.}). (\sou{kemio)}. Lumar en loko obskura. (ex.:}
\l{\cel{3}fosfo). - DEFIRS.}
\l{}
\l{\sou{fosilo}. Qua permanis enfosigita en la strati sedimenta anciena}
\l{\cel{3}di la Terkortico. - DEFIS.}
\l{}
\l{\sou{fosto}. I. Peco de karpenturo, pozita vertikale. - II. Longa}
\l{\cel{3}ligno-peco stekita en sulo. - D.}
\l{}
\l{\sou{foto}. (\sou{fiziko)}. Radiado de ula korpi, qua videbligas la ob-}
\l{\cel{3}jekti.}
\l{}
\l{\sou{fotelo}. Granda stulo, kun dors-apogilo e brakii. - F. Pol.}
\l{}
\l{\sou{fotofono}. Aparato de ube la son-ondi, recevite sur spegulo}
\l{\cel{3}vibranta qua transformas li ad lum-ondi, departas por re-}
\l{\cel{3}produktar su kelka-diste. - DEFIRS.}
\l{}
\l{\sou{fotogena}.\cel{2}Dicesas pri irga parto di organismo qua produktas}
\l{\cel{3}ed emisas lumo-radii fiziologiala.}
\l{}
\l{\sou{fotograbar}. (\sou{trans}.) Aplikar la fotografado e la grabado per}
\l{\cel{3}vitro-plako lumo-impresebligita. - DEFIRS.}
\l{}
\l{\sou{fotografar.}\cel{1}(\sou{trans.})\cel{2}Fixigar la imajo de korpi, obtenita per}
\l{\cel{3}kamero obskura, sur vitroplako (o filmo) garnisita ye sub-}
\l{\cel{3}stanco impresebla da lumo. - DEFIRS.}
\l{}
\l{\sou{fotokemio}. La cienco di la reakti kemiala quin produktas la}
\l{\cel{3}lumo-radii.}
\l{}
\l{\sou{fotolitografar}. (\sou{trans.)}\cel{2}Aplikar la fotografado a la lito-}
\l{\cel{3}grafado per fotografuro dekalquebla sur petro.}
\l{}
\l{\sou{fotometro}. Instrumento por mezurar la intenseso-grado di lumo.}
\l{\cel{3}- DEFIRS.}
\l{}
\l{\sou{fotometrio}. Parto di fiziko, qua traktas la mezuro di la in-}
\l{\cel{3}tenseso-gradi lumala.}
\l{}
\l{\sou{fotosfero}. (\sou{astron.})\cel{2}La atmosfero lumoza qua cirkondas Suno.}
\l{\cel{3}- DEFIRS.}
\l{}
\l{\sou{fotosintezo.}- (\sou{bot.)}\cel{1}Sintezo qua eventas en la folii di la}
\l{\cel{3}planti verda, influe da la lumo-radii.}
\l{}
\l{\sou{fototaktismo.}\cel{1}Movo (influe da lumo-radii) da korpi moviva}
\l{\cel{3}(zoospori, e c.)}
\l{\sou{fototerapio.}\cel{1}Uzo di la lumo-radii kom agento terapiala.-DEF.}
}
